\documentclass[12pt]{article}
\usepackage[utf8]{inputenc}
\usepackage{amsmath}
\usepackage{graphicx}
\usepackage{natbib}
\usepackage{hyperref}
\usepackage{lipsum} % For dummy text, you can remove this in your final document

\title{Impact of Urbanization on Biodiversity: Habitat Loss, Species Extinction, Ecosystem Services, Mitigation Strategies, and Recommendations}
\author{Your Name}
\date{\today}

\begin{document}

\maketitle

\begin{abstract}
Urbanization is a significant driver of biodiversity loss worldwide, impacting habitats, species diversity, and ecosystem services. This paper reviews the effects of urbanization on biodiversity, analyzes case studies from different regions, proposes mitigation strategies, and provides recommendations for sustainable urban development.
\end{abstract}

\section{Introduction}
Urbanization has rapidly transformed landscapes globally, leading to substantial changes in natural habitats and ecosystems. The expansion of cities and infrastructure development increasingly encroaches upon biodiversity-rich areas, affecting species survival, ecosystem functions, and the provision of essential services. This paper examines the multifaceted impacts of urbanization on biodiversity, focusing on habitat loss, species extinction risks, and alterations in ecosystem services.

\section{Impact Analysis}
\subsection{Habitat Loss}
Urban development results in the fragmentation and degradation of natural habitats, reducing available space for wildlife and altering landscape connectivity \citep{mcdonnell2015ecology, forman2014urban}.

\subsection{Species Extinction}
The conversion of natural habitats into urban landscapes increases the risk of species extinction due to habitat destruction, isolation, and human-wildlife conflicts \citep{loss2018biodiversity, gaston2006urban}.

\subsection{Ecosystem Services}
Urbanization alters ecosystem functions such as water purification, climate regulation, and pollination services, impacting human well-being and urban resilience \citep{elmqvist2013urban, haase2014challenges}.

\section{Case Studies}
\subsection{North America}
Case studies in North American cities highlight biodiversity declines, including loss of native species and disruptions in ecological processes \citep{mcdonald2008urbanization, aronson2014biodiversity}.

\subsection{Asia}
Urban expansion in Asian cities contributes to biodiversity loss through habitat destruction, pollution, and land use changes \citep{seto2012global, gopal2017urban}.

\section{Mitigation Strategies}
\subsection{Green Infrastructure}
Implementing green spaces, parks, and urban forests to preserve biodiversity, promote habitat connectivity, and enhance urban ecosystem services \citep{marusenko2019green, miller2019urban}.

\subsection{Land Use Planning}
Integrating biodiversity conservation into urban planning processes, including zoning regulations, habitat restoration programs, and sustainable development practices \citep{alberti2017urban, collins2011biodiversity}.

\section{Recommendations}
\subsection{Policy Interventions}
Enacting and enforcing policies that prioritize biodiversity conservation in urban areas, incentivize green development, and mitigate urban sprawl \citep{marull2015social, mcdonnell2016biodiversity}.

\subsection{Community Engagement}
Engaging local communities in biodiversity monitoring, habitat restoration initiatives, and promoting awareness of urban biodiversity values \citep{pereira2019engaging, shwartz2016urban}.

\section*{References}
\bibliographystyle{plainnat}
\bibliography{prompt13_35}

\end{document}
